\section{Details of the Examples}
\label{app:examples}

\subsection{The Lambda Class of Example 2}
\label{app:lambda}
\texttt{javac} compiles the body of the lambda \Q@s->s@ of Example~2 into a private static method \Q@lambda$main$1(String)@ of the class \Q@Broken@.
At run time, the first evaluation of the lambda \Q@s->s@ calls \Q@LambdaMetafactory@, which generates the lambda class.
The method \Q@apply(Object)@ of the lambda class casts its argument to \Q@String@, and then calls \Q@lambda$main$1@ with the result of that cast.

The stack trace of the \Q@ClassCastException@ does not show the lambda class.
\Q@LambdaMetafactory@ defines the lambda class as a hidden class, and the JVM omits the frames of the methods of hidden classes from stack traces, unless the JVM runs with the diagnostic option \Q@-XX:+ShowHiddenFrames@.
So the topmost frame of the stack trace is the frame of \Q@m@, at Line~\ref{line:ex2apply} of Example~2.
The bytecode of \Q@m@ at Line~\ref{line:ex2apply} does contain a \Q@checkcast@ instruction to \Q@String@, on the result of \Q@f.apply@, but that instruction is never reached.

\subsection{Example 2 with a \Q@switch@}
\label{app:switch}
This is the full program of Example~2 with a \Q@switch@ instead of the \Q@if@.
javac accepts this program with no warning, and this program fails as Example~2 fails.

\noindent\begin{minipage}{\linewidth}
\begin{lstlisting}[numbers=left,xleftmargin=2em]
// Example 2 with a switch
public class Broken {
  static <X> Object m(Object other, final X x, Function<X,String> f){
    class MyFoo{ public X x(){ return x; } }
    switch(other){
      case MyFoo local: // no cast, no warnings
        System.out.println("old was "+f.apply(local.x())); //ClassCastException
        break;
      default:
        break;
    }
    return new MyFoo();
  }
  public static void main(String[] a) {
    Object r1= Broken.<Integer>m("anything",5,i->""+i);
    Broken.<String>m(r1,"Hi",s->s);
  }
}
\end{lstlisting}
\end{minipage}
